\section*{Einleitung}
Im Rahmen des Moduls STEP (Einführung in die Informatik/Wirtschaftsinformatik) hatte jedes Team die Aufgabe, auf Basis von AIS-Daten (Automatisches Identifikationssystem) ein Projekt durchzuführen. Bei AIS handelt es sich um ein Funksystem im Ultrakurzwellenbereich (UKW), welches automatiert Schiffinformationen für weitere Schiffe und Behörden zur Verfügung stellt ~\autocite[vgl.][S.32]{fischer2016}. Die wichtigsten Funktionen sind dabei Kollisionsverhinderung, Erhalt von Schiffsinformationen und Ladung sowie eine Verbindungstelle für die Regelung des Schiffsverkehrs vom Land ~\autocite[vgl.][S.32]{fischer2016}. 
